\begin{figure}[htbp]
\centering
\begin{tikzpicture}[font=\small,>=Stealth]
\node[font=\small\bfseries] at (0,1.6) {Trace-sensitive quotient};
\node[font=\small\bfseries] at (5.4,1.6) {Observable quotient};
\draw[pathblue,thick,rounded corners] (-1.65,-0.45) rectangle (-0.35,0.75);
\draw[pathblue,thick,rounded corners] (0.35,-0.45) rectangle (1.65,0.75);
\fill[pathblue] (-1,0.3) circle (2pt);
\fill[pathblue] (1,0.3) circle (2pt);
\node[below] at (-1,0.15) {$[u_e]$};
\node[below] at (1,0.15) {$[u_f]$};
\draw[pathorange,thick,rounded corners] (3.7,-0.45) rectangle (7.1,0.75);
\fill[pathorange] (4.4,0.3) circle (2pt);
\fill[pathorange] (6.4,0.3) circle (2pt);
\node[below] at (4.4,0.15) {$[u_e]$};
\node[below] at (6.4,0.15) {$[u_f]$};
\draw[->,thick] (1.95,0.3) -- (3.35,0.3) node[midway,above] {$J$};
\node[align=center] at (0,-1) {Distinct classes;\\separated by an open set};
\node[align=center] at (5.4,-1) {Distinct classes;\\identical open neighborhoods};
\end{tikzpicture}
\caption{The two quotients identify the same arrows and differ only in
their topology.
In Example~\ref{ex:quotient-topology-separation}, both steps realize the
same nonconstant circle loop, but their signed-count values differ. The map $J$
preserves both scoped classes.  In the observable quotient every open set
containing either class contains both.  Boxes illustrate separation properties
only; the shared box is not asserted to be an open two-point subset.}
\label{fig:same-geometry}
\end{figure}
